% 中文正文片段；主文件加载 amsmath, amssymb。
% 历史批次 20260922T023439-35a9711f；报告 20260922T032418-3cd92bd4。
\subsection{历史定常外力方腔的参考解时间误差}

本研究整理固定空间网格上的普通 IMEX--SDIRK2 与增量
SDIRK2--mr-ccSAV 时间误差。区域为单位正方形 $\Omega=(0,1)^2$，
边界为齐次无滑移，采用 $128\times128$ 均匀交错 MAC 网格，
运动黏性系数 $\nu=0.01$，零速度初值，SAV 阻尼参数 $\gamma=1$。
必须区分历史记录与当前生产配置：本节所用运行配置明确记录
\begin{equation}
 f(x,y)=(0,\sin x),
\end{equation}
其中 $x$ 是物理坐标。本节不使用当前配置的
$f=(0,\sin(4\pi x))$，也不将当前外力下的定常窗口结论转用于本批次。
两方法在相同网格与物理设置上推进至 $T=2$，在 $t=1,1.5,2$ 比较速度。

被测步长为 $\tau_k=0.1\,2^{-k}$，$k=1,\ldots,8$，
共十六个被测运行。共同参考采用普通 SDIRK2 的 $k=12$ 结果，
再以 $k=13$ 结果检查参考敏感性。设数值速度与共同参考之差为
$e_h$，误差指标为
\begin{equation}
 E_{L^2}=\sqrt{(e_h,e_h)_h},\qquad
 E_{H^1}=\sqrt{(e_h,K_he_h)_h},
\end{equation}
其中 $(\cdot,\cdot)_h$ 为 MAC 速度内积，$K_h$ 为正的离散黏性算子，
第二项为速度的离散 $H^1$ 半范数误差，不是压力误差。
相邻步长层的有效观察阶为 $p_k=\log_2(E_{k-1}/E_k)$。

既有误差表显示，在 $T=2$，普通格式的 $L^2$ 误差由
$k=1$ 的约 $3.84\times10^{-6}$ 减至 $k=8$ 的约
$2.34\times10^{-10}$；最细层的 $L^2$ 观察阶约为 $2.001$，
$H^1$ 半范数观察阶约为 $1.995$。相应 SAV 格式在本批次给出接近的误差，
$H^1$ 半范数观察阶逐步趋近二阶，而不是要求所有粗步长层均已有形式阶。
这些是所选离散系统和观测窗口中的数值观察，不表示两种格式在一般问题上等价。

两个参考解的差除以每一观测时刻和范数下最小有效被测误差，
六项保存检查均不超过 $5\%$，最大约为 $3.64\%$。
该标准仅衡量所用参考细化对的敏感性，不构成严格参考误差界。
展示从固定的历史误差表重绘，并核对全部十八份输入的配置、来源和结果摘要，
不重新推进或独立重算场误差。单一固定空间网格上的时间比较不能替代
空间收敛验证或无滑移偏微分方程的高阶收敛证明。
